% Page budget: 1.55 pages | Owns: negation, OBC derivation, additional-state scope, implementation, and true-parameter condition; per README ownership also the model M_OBC.
% WRITING GUIDE
% Purpose: Explain negation, derive the final state-level orthogonal-by-construction map, and delimit its interpretability guarantee.
% Sources: README "Source map per subsection", rows IV-A, IV-B and IV-C; mandatory papers and the reference gate as in the README.
% Research-plan anchor: Preserve Aspect 3's goal of protecting physical interpretability; state clearly how the final construction differs from the proposed regularisation.
% Current decisions: Reduced identifiable coordinates, data-derived reference tuples, a tangent-only one-step physical-state projection, per-epoch refresh, and no affine comparison arm.
% Choices to justify: Protected parameter coordinates, reference-set construction, projected rows, rank tolerance, exclusion of the affine offset, refresh frequency, and treatment of additional states.
% Cautions: docs/orthogonal-projection-plan.md describes an older soft-penalty route; one-step state orthogonality is not trajectory-output orthogonality or guaranteed true recovery.
% Planned figures: New geometric projection figure with a physical and additional-state partition inset; show the remaining uniqueness limitation as clearly as the protected direction.
% Section is finished when: The protected space, coordinate frame, gradient treatment, added-state scope, and true-recovery condition are explicit.
%
% PROPOSED MUST ESTABLISH (writer, sandbox cycle 05, 2026-10-09; NOT approved by Dirk; inferred from the README
%   ownership row IV, source-map rows IV-A to IV-C, DOC/THESIS-RESULTS.md step 5, the Meeting-audit OBC theme and the
%   Section III handover "Section IV extends the method so that the estimate is unique")
% Contribution delivered: third contribution of Section I (state-level orthogonality for the identifiable
%   physical-parameter tangent space and its true-parameter recovery condition).
% IV-A 1. Negation (Kessels Sec. 5.3.1.3; Gyorok 2026 Ex. 2) in its gantry form: the transition sensitivity Phi with
%      respect to the training coordinates xi gives a first-order cancellation; its columns act through
%      qdd, qd, q with M(Y)^-1 (own, checked by differentiating eq:eom).
% IV-B 2. Adopted: Gyorok 2026 Eqs. 8, 10, Lemma 3, Eq. 13. Adapted: regressor = Jacobian of the normalised RK4
%      transition at xi_e (Taylor precedent Gyorok 2025 Sec. 4), on data-derived tuples (q = P^-T y, fourth-order
%      difference, xbar = 0, recorded input), refreshed per epoch, no offset column. Own: per-objective coefficient
%      inside the closed-loop rollout, model M_OBC with its full identification problem (README Math standard 3).
% IV-C 3. Proposition (own): state-level counterpart of Gyorok 2026 Thm. 7 / Eq. 21: unique estimate for every
%      network, true one iff Phi' delta = 0. Scope: one-step surrogate, xbar = 0 slice, g_aug uncorrected,
%      normalised inner product, Thms. 16, 18 do not transfer. Instance (T_AF friction) left to Section V.
% Left out (reader test or ownership): hand-written forward sensitivity (D-194: implementation of the same
%   derivative), SVD and rank tolerance (header list item, fails reader test), per-objective hoist (D-195),
%   validation-sync mechanics beyond one clause (D-198), overlap measures (not a THESIS-RESULTS metric), affine arm
%   (not in the campaign, THESIS-RESULTS), noise check of the basis (D-232 result; observed quantity).
% Existing todos resolved in this draft: "Source" (opening now cites Gyorok 2026 Abstract and Sec. 6); "Subsection
%   title improve" (new title); "From Section III ... owns the term negation" (IV-A opens from it); affine-arm
%   wording and its \note (THESIS-RESULTS has no affine arm); refs.bib \note (entries added); "one-step vs
%   trajectory penalty" (D-190 Ruled out 1: trajectory method rejected by the user); "orthogonal multisine"
%   (THESIS-RESULTS Sec. 6: T_OA not run; no constructive remedy is claimed). Kept: Gamma exclusion todo
%   (supervisor confirmation).
% Notation: Phi, xi_e, theta_aux, N_r, bold stacks, delta, xi^star, xi^circ are new; checked against
%   util/format.tex and sections 00 to 08, 99. Later sections use J and Delta^* (todo in IV-C).
% Kept from the reset header notes: Figure plan (geometric decomposition with x_a inset) not drawn: figures/ is
%   generated, no route in this session.
\section{Interpretability by Orthogonal Construction}\label{sec:obc}
\ifdraft
\begin{itemize}
\item Given the non-unique split of Section~III-A, the goal is a parametrisation in which the estimate of $\theta_{\mathrm{base}}$ is unique; the section names its three subsections (style profile)
\item Contribution share: the third contribution of Section~I, state-level orthogonality for the identifiable tangent space and its true-parameter recovery condition (Introduction)
\end{itemize}
\fi

Section~\ref{sec:aug_structure} left the split between the baseline and the
network non-unique.
Our goal is a parametrisation of the augmented model in which the estimate of
$\theta_{\mathrm{base}}$ is unique.
To this end, Section~\ref{sec:obc_negation} states how the network can
negate the gantry baseline.
Section~\ref{sec:obc_map} removes the parameter directions from the network
write by construction and defines $\mathcal M_{\mathrm{OBC}}$.
Section~\ref{sec:obc_scope} states what the construction guarantees and when
it recovers the true parameters.
Together they construct state-level orthogonality for the identifiable
physical-parameter tangent space and quantify its true-parameter recovery
condition, the third contribution of Section~I.

\subsection{Negation in the gantry transition}\label{sec:obc_negation}
\ifdraft
\begin{itemize}
\item Joint estimation lets the network negate baseline terms (Kessels Sec.~5.3.1.3); for linear models every pair with the same sum fits (Gy\"or\"ok 2026 Ex.~2)
\item On the gantry the exchange is local: the transition sensitivity $\Phi$ with respect to $\xi$ gives a first-order cancellation (this work)
\item $\Phi$ has the range of the $\theta_{\mathrm{base}}$ sensitivity, whose columns are RK4 propagations of $\partial\ddot q/\partial\theta_{\mathrm{base},j}$; a negating write is an added inertial, damping or stiffness force scaled by $M(Y)^{-1}$ (this work)
\end{itemize}
\fi

With $\theta_{\mathrm{base}}$ and $\theta_{\mathrm{aug}}$ estimated jointly,
the network can represent dynamics that the baseline already represents.
Kessels calls learned terms that (partially) cancel the first-principles terms
negation~\cite[Sec.~5.3.1.3]{kessels2025ai}.
For a baseline and a network layer that are both linear, every parameter pair
with the same sum minimises the cost~\cite[Ex.~2]{gyorok2026obc}.

On the gantry, the exchange is local in the training coordinates $\xi$ of
\eqref{eq:mdelta_map}.
Its directions are the columns of the sensitivity of the normalised transition
of \eqref{eq:aug_norm},
\begin{equation}
 \Phi(\xi,\tilde x^{\mathrm n},u^{\mathrm n})
 =\frac{\partial f^{\mathrm n}_{\mathrm{base}}
   \bigl(\theta_{\mathrm{base}}(\xi),\tilde x^{\mathrm n},u^{\mathrm n}\bigr)}
   {\partial\xi},
 \label{eq:obc_sens}
\end{equation}
where $\Phi:\R^{10}\times\R^6\times\R^3\to\R^{6\times10}$ and
$\theta_{\mathrm{base}}(\xi)$ is the map \eqref{eq:mdelta_map}.
Negation is then precise for the parameter directions: for every
displacement $\Delta\xi\in\R^{10}$,
\begin{equation}
\begin{aligned}
 &f^{\mathrm n}_{\mathrm{base}}\bigl(\theta_{\mathrm{base}}(\xi+\Delta\xi),\tilde x^{\mathrm n},u^{\mathrm n}\bigr)
 -\Phi(\xi,\tilde x^{\mathrm n},u^{\mathrm n})\,\Delta\xi\\
 &\quad=f^{\mathrm n}_{\mathrm{base}}\bigl(\theta_{\mathrm{base}}(\xi),\tilde x^{\mathrm n},u^{\mathrm n}\bigr)
 +O(\norm{\Delta\xi}^2).
\end{aligned}
\label{eq:obc_negation}
\end{equation}
A network that subtracts $\Phi\,\Delta\xi$ from its normalised write
$T_xf_{\mathrm{aug}}$ thus leaves the transition unchanged to first order
while $\xi$ moves by $\Delta\xi$.

The columns of $\Phi$ have a physical reading.
The Jacobian of \eqref{eq:mdelta_map} is invertible, so $\Phi$ has the range
of the sensitivity with respect to $\theta_{\mathrm{base}}$.
RK4 propagates into that sensitivity the derivative of \eqref{eq:eom} at fixed
$q$, $\dot q$ and $u$,
\begin{equation}
 \frac{\partial\ddot q}{\partial\theta_{\mathrm{base},j}}
 =-M(Y)^{-1}\Bigl(\frac{\partial M(Y)}{\partial\theta_{\mathrm{base},j}}\ddot q
 +\frac{\partial C}{\partial\theta_{\mathrm{base},j}}\dot q
 +\frac{\partial K}{\partial\theta_{\mathrm{base},j}}q\Bigr),
 \label{eq:obc_gantry_sens}
\end{equation}
where $j$ indexes the ten entries of \eqref{eq:phi}.
A negating write therefore acts as an added inertial, damping or stiffness
force, scaled by $M(Y)^{-1}$.

\subsection{Orthogonal-by-construction correction}\label{sec:obc_map}
\ifdraft
\begin{itemize}
\item Gy\"or\"ok et al.\ remove negation by construction for IO models linear in the parameters, without the penalty weight of their 2025 method; state space is future work there (Gy\"or\"ok 2026 Abstract, Sec.~6; Gy\"or\"ok 2025 Eq.~(13), Remark~1)
\item Adapted: regressor = $\Phi$ at an expansion point $\xi_{\mathrm e}$, as the Taylor expansion of Gy\"or\"ok 2025 Sec.~4
\item Reference tuples from every training sample: $P^{-\top}y$, fourth-order difference, $\bar x=0$, recorded input, normalised (code \texttt{build\_reference\_set})
\item Coefficient, corrected stack and orthogonality on the stack, rank 10 required (Gy\"or\"ok 2026 Eqs.~(8), (10), Lemma~3, Assumption~1)
\item $\mathcal M_{\mathrm{OBC}}$: identification problem of Section~III-C with the corrected transition, in one display (this work)
\item Lifecycle: $\theta_{\mathrm{aux}}$ per evaluation of $V$, $\xi_{\mathrm e}$ per epoch, fixed $\theta_{\mathrm{aux}}$ for prediction (Gy\"or\"ok 2026 Eq.~(13)); start, PS2, optimiser and selection as $\mathcal M_{\mathrm{PS2}}$ (D-192, D-198)
\item Realisation choices in one paragraph: ten coordinates, fixed tuples (causality), $\bar x=0$ (?), no offset column (unlike Gy\"or\"ok 2025 Eq.~(18)), refresh per epoch (D-190, D-192)
\end{itemize}
\fi

Gy\"or\"ok et al.\ remove negation by construction for input-output models
that are linear in their parameters~\cite[Abstract, Sec.~2]{gyorok2026obc}.
They subtract from the network output its least-squares fit on the stacked
baseline regressor~\cite[Eqs.~(8), (10)]{gyorok2026obc}.
Unlike the projection penalty of~\cite[Eq.~(13)]{gyorok2025l4dc}, whose weight
trades orthogonality against fit~\cite[Remark~1]{gyorok2025l4dc}, this
correction has no weight to tune.
Gy\"or\"ok et al.\ leave the extension to state-space models without
full-state measurement open~\cite[Sec.~6]{gyorok2026obc}.
We make it for the gantry.
The gantry transition is rational in $\theta_{\mathrm{base}}$ through
$M(Y)^{-1}$, so we take $\Phi$ at an expansion point $\xi_{\mathrm e}$ as the
regressor, as the Taylor expansion of~\cite[Sec.~4]{gyorok2025l4dc} does for
the penalty.

The regressor is evaluated on states, while only positions are measured.
We therefore form reference tuples from every training sample with a full
difference stencil, $i=1,\dots,N_{\mathrm r}$,
\begin{equation}
 \tilde x_i=\col\bigl(P^{-\top}y_i,\;\dot q_i\bigr),\qquad
 \bar x_i=0,\qquad
 u_{\mathrm{act},i},
 \label{eq:obc_tuples}
\end{equation}
where $y_i\in\R^3$ and $u_{\mathrm{act},i}\in\R^3$ are the recorded output and
input, $\dot q_i\in\R^3$ is the fourth-order central difference of
$P^{-\top}y$ within the record of sample $i$, and
$\tilde x^{\mathrm n}_i$, $u^{\mathrm n}_i$ are the tuples normalised as in
\eqref{eq:aug_norm}.

On these tuples, the construction of~\cite[Eqs.~(8), (10), Lemma~3]{gyorok2026obc}
becomes
\begin{equation}
\begin{aligned}
 \Phi(\xi_{\mathrm e})
 &=\col\bigl(\Phi(\xi_{\mathrm e},\tilde x^{\mathrm n}_i,u^{\mathrm n}_i)\bigr)_{i=1}^{N_{\mathrm r}},\\
 \mathbf f_{\mathrm{aug}}(\theta_{\mathrm{aug}})
 &=\col\bigl(T_xf_{\mathrm{aug}}(\theta_{\mathrm{aug}},\col(\tilde x_i,0),Pu_{\mathrm{act},i})\bigr)_{i=1}^{N_{\mathrm r}},\\
 \theta_{\mathrm{aux}}&=\Phi(\xi_{\mathrm e})^{+}\,\mathbf f_{\mathrm{aug}}(\theta_{\mathrm{aug}}),\\
 0&=\Phi(\xi_{\mathrm e})^\top
 \bigl(\mathbf f_{\mathrm{aug}}(\theta_{\mathrm{aug}})-\Phi(\xi_{\mathrm e})\,\theta_{\mathrm{aux}}\bigr),
\end{aligned}
\label{eq:obc_construction}
\end{equation}
where $\Phi(\xi_{\mathrm e})\in\R^{6N_{\mathrm r}\times10}$ is the stacked
sensitivity, $\mathbf f_{\mathrm{aug}}(\theta_{\mathrm{aug}})\in\R^{6N_{\mathrm r}}$
the stacked normalised network write, $\theta_{\mathrm{aux}}\in\R^{10}$ the
auxiliary coefficient, and $(\cdot)^+$ the pseudoinverse.
As in~\cite[Assumption~1]{gyorok2026obc}, $\Phi(\xi_{\mathrm e})$ must have full
column rank.
\todo{Missing: where the rank of $\Phi(\xi_{\mathrm e})$ on the thesis records is reported (README ownership: observed quantities belong to Results or the appendix, and THESIS-RESULTS lists no rank metric). Candidate: one appendix line from the run logs, which print rank and condition number at every refresh (\texttt{OBCLifecycle.refresh}; D-232 smoke: rank 10/10, condition number 656).}

In the rollout, the fitted directions are subtracted at the current state.
The identification problem of $\mathcal M_{\mathrm{OBC}}$ is that of
Section~\ref{sec:aug_closed_loop} with this corrected transition,
\begin{equation}
\begin{aligned}
 &\min_{\vartheta}\;V(\vartheta)\quad\text{subject to}\\
 &\tilde x_{k+1}=f_{\mathrm{base}}(\theta_{\mathrm{base}}(\xi),\tilde x_k,u_k)
   +f_{\mathrm{aug}}(\theta_{\mathrm{aug}},\hat x_k,u_k)\\
 &\qquad\quad-T_x^{-1}\Phi(\xi_{\mathrm e},\tilde x^{\mathrm n}_k,u^{\mathrm n}_k)\,\theta_{\mathrm{aux}},\\
 &\bar x_{k+1}=g_{\mathrm{aug}}(\theta_{\mathrm{aug}},\hat x_k,u_k),\qquad
   \hat y_k=h_{\mathrm{base}}(\tilde x_k),\\
 &\theta_{\mathrm{aux}}=\Phi(\xi_{\mathrm e})^{+}\,\mathbf f_{\mathrm{aug}}(\theta_{\mathrm{aug}}),\\
 &\col\bigl(T_x(\tilde x_{\tau}-\mu_x),\bar x_{\tau}\bigr)
   =\psi(\theta_{\mathrm{enc}},\mathbf y_\tau,\mathbf u_\tau),\\
 &x^{c}_{k+1}=A_{\mathrm{fb}}x^{c}_k+B_{\mathrm{fb}}(y_k-\hat y_k),\qquad x^{c}_\tau=0,\\
 &u_k=P\bigl(u_{\mathrm{act},k}+C_{\mathrm{fb}}x^{c}_k+D_{\mathrm{fb}}(y_k-\hat y_k)\bigr),
\end{aligned}
\label{eq:obc_model}
\end{equation}
where $\vartheta=(\xi,\theta_{\mathrm{aug}},\theta_{\mathrm{enc}})$, $V$ is
\eqref{eq:aug_loss}, $k$ runs over every window of length $n_f$ that starts
at $\tau\in\mathcal T$, and $\tilde x^{\mathrm n}_k$, $u^{\mathrm n}_k$ are
$\tilde x_k$ and the actuator force $P^{-1}u_k$ normalised as in
\eqref{eq:aug_norm}.
We call this model $\mathcal M_{\mathrm{OBC}}$.
\todo{Missing: model notation. Candidate: $\mathcal M_{\mathrm{OBC}}$, since $M$ is the inertia matrix (README open conflict on model names, as for $\mathcal M_{\mathrm{U}}$ and $\mathcal M_{\mathrm{PS2}}$ in Section~\ref{sec:augmentation}).}

The coefficient and the expansion point are updated at different rates.
$\theta_{\mathrm{aux}}$ is recomputed from the current network at every
evaluation of $V$, so the gradient of $V$ passes through the correction.
$\xi_{\mathrm e}$ is set to the current $\xi$ at every epoch boundary and at
the selected checkpoint.
$\Phi(\xi_{\mathrm e})$ is held fixed in between.
Validation and test rollouts use $\theta_{\mathrm{aux}}$ computed once from the
current network, as in~\cite[Eq.~(13)]{gyorok2026obc}.

$\mathcal M_{\mathrm{OBC}}$ starts and is trained as $\mathcal M_{\mathrm{PS2}}$.
At the zero start \eqref{eq:aug_zero_init}, $\mathbf f_{\mathrm{aug}}=0$, so
$\theta_{\mathrm{aux}}=0$ and the model starts from the baseline.
It is then initialised by PS2, optimised and selected as in
Sections~\ref{sec:aug_ps2} and~\ref{sec:aug_closed_loop}.

The realisation choices serve a unique $\xi$ at an affordable cost.
$\xi$ parameterises only the ten identifiable combinations, so $\Phi$ can
have full column rank.
In the fourteen raw parameters, the four lost directions of
Section~\ref{sec:params} would be null directions of the stack.
The tuples depend only on the records.
$\theta_{\mathrm{aux}}$ therefore depends only on $\theta_{\mathrm{aug}}$ and
exists before the rollout it corrects, whereas a fit on the rollout states
would make each step depend on later states.
Such an auxiliary evaluation set is permitted by~\cite[Sec.~3.2]{gyorok2026obc}.
Unlike~\cite[Eq.~(18)]{gyorok2025l4dc}, the stack omits the offset
$f^{\mathrm n}_{\mathrm{base}}(\theta_{\mathrm{base}}(\xi_{\mathrm e}),\cdot)$,
because the uniqueness of $\xi$ depends on $\Phi$ alone
(Section~\ref{sec:obc_scope}).
The one-step offset is dominated by the propagated state, so protecting it
would also remove network directions that cause no parameter ambiguity.
Gy\"or\"ok et al.\ update the expansion point at every evaluation of the
cost or fix it at the nominal parameters~\cite[Sec.~4]{gyorok2025l4dc}.
Updating at every evaluation rebuilds $\Phi$ on all tuples each time.
The range of $\Phi$ turns slowly with $\xi$, so we refresh it once per epoch.
A fixed nominal point does not apply, because training may start detuned
(Section~\ref{sec:setup}).
\todo{Missing: reason for evaluating the tuples at $\bar x=0$. No source states one (D-192 item~2 and the one-step implementation plan only state the choice). Candidate: the added states are not measured and have no physical coordinates (Section~\ref{sec:aug_structure}), so the records define no value for them, and $\bar x=0$ is their value at the zero start (own reasoning).}
\todo{Gamma excluded: the protected set is the range of Phi (parameter uniqueness, \cite[Assumption~1, Thm.~7]{gyorok2026obc}), not the range of [Phi Gamma] (output overlap, \cite[Eq.~(18)]{gyorok2025l4dc}). Confirm with supervisors, since D-190 records their condition as non-overlap with the baseline output.}

\subsection{Scope and true-parameter recovery}\label{sec:obc_scope}
\ifdraft
\begin{itemize}
\item Proposition: at a stationary point of the one-step criterion on the tuples, $\xi^\circ-\xi^\star=\Phi^+\bm\delta$ for every network, and $\xi^\circ=\xi^\star$ iff $\Phi^\top\bm\delta=0$ (this work; counterpart of Gy\"or\"ok 2026 Thm.~7, Eq.~(21))
\item Without the correction the estimate moves with the network (Gy\"or\"ok 2026 Sec.~4.2, Eq.~(25)); with it the estimate is unique, the requirement of Section~III-A
\item Uniqueness is not recovery: Condition~4 in state form; designed excitation can meet it (Gy\"or\"ok 2026 Cond.~4, Remark~5); whether the simulated truth meets it belongs to Section~V (?)
\item Limits: one-step surrogate of $V$; tuples at $\bar x=0$, so the added-state route of $f_{\mathrm{aug}}$ is unconstrained and $g_{\mathrm{aug}}$ uncorrected; normalised inner product; Thms.~16, 18 do not transfer (D-192 Constrains; AUGMENTATION-INITIALIZATION 12.2; THESIS-RESULTS step~5)
\end{itemize}
\fi

The construction makes the parameter estimate independent of the network
when the training criterion is the one-step error on the tuples.
Let $\mathbf x^+\in\R^{6N_{\mathrm r}}$ stack the normalised recorded
successors $\tilde x^{\mathrm n}_{i+1}$ of the tuples, formed with the same
stencil, and let
$\mathbf f_{\mathrm{base}}(\xi)=\col\bigl(f^{\mathrm n}_{\mathrm{base}}(\theta_{\mathrm{base}}(\xi),\tilde x^{\mathrm n}_i,u^{\mathrm n}_i)\bigr)_{i=1}^{N_{\mathrm r}}$.
With $\xi^\star$ the coordinates of the true parameters, the one-step
discrepancy on the tuples is
\begin{equation}
 \bm\delta=\mathbf x^+-\mathbf f_{\mathrm{base}}(\xi^\star),
 \label{eq:obc_delta}
\end{equation}
where $\bm\delta\in\R^{6N_{\mathrm r}}$ contains every effect the baseline
omits, measurement noise included.

\textit{Proposition 1:} Suppose (i) $\mathbf f_{\mathrm{base}}$ is affine in
$\xi$ on a set that contains $\xi^\star$, $\xi_{\mathrm e}$ and $\xi^\circ$,
so that $\Phi=\Phi(\xi_{\mathrm e})$ is its Jacobian there; (ii) $\Phi$ has
full column rank; and (iii) for a given $\theta_{\mathrm{aug}}$, $\xi^\circ$ is
a stationary point in $\xi$ of
$\norm{\mathbf x^+-\mathbf f_{\mathrm{base}}(\xi)
-\bigl(\mathbf f_{\mathrm{aug}}-\Phi\,\theta_{\mathrm{aux}}\bigr)}_2^2$,
with $\theta_{\mathrm{aux}}$ from \eqref{eq:obc_construction}.
Then $\xi^\circ-\xi^\star=\Phi^+\bm\delta$, whatever $\theta_{\mathrm{aug}}$,
and $\xi^\circ=\xi^\star$ if and only if $\Phi^\top\bm\delta=0$.

\begin{IEEEproof}
By (i), $\mathbf f_{\mathrm{base}}(\xi^\circ)=\mathbf f_{\mathrm{base}}(\xi^\star)
+\Phi(\xi^\circ-\xi^\star)$.
Stationarity in $\xi$ requires $\Phi^\top r=0$ for the residual
$r=\bm\delta-\Phi(\xi^\circ-\xi^\star)
-(\mathbf f_{\mathrm{aug}}-\Phi\,\theta_{\mathrm{aux}})$.
The last line of \eqref{eq:obc_construction} removes the network term, so
$\Phi^\top\Phi(\xi^\circ-\xi^\star)=\Phi^\top\bm\delta$, and (ii) gives both
claims.
\end{IEEEproof}

Proposition~1 is our state-level counterpart
of~\cite[Thm.~7, Eq.~(21)]{gyorok2026obc}.
Without the correction, the same steps give
$\xi^\circ-\xi^\star=\Phi^+(\bm\delta-\mathbf f_{\mathrm{aug}})$, so the
estimate moves with the network, as in~\cite[Sec.~4.2]{gyorok2026obc}.
With the correction, the estimate is unique, which is the requirement of
Section~\ref{sec:aug_structure}.
Because $\xi\mapsto\theta_{\mathrm{base}}(\xi)$ is injective, the same holds
for $\theta_{\mathrm{base}}$.
\todo{Notation conflict: Sections~\ref{sec:results} and VII and the appendix bullets write $J$ and $\Delta^*$ for $\Phi$ and $\bm\delta$, but $\Delta$ is the scheduling block of Section~\ref{sec:lfr} and $J$ the inertias $J_b$, $J_h$, $J_{\mathrm{eff}}$. Candidate: use $\Phi$ and $\bm\delta$ there. The Section~\ref{sec:results} bullet also cites an overlap display \texttt{eq:obc\_overlap}, which this section no longer defines because THESIS-RESULTS step~5 lists no overlap metric.}

Uniqueness is not recovery.
The true coordinates follow only if the discrepancy is orthogonal to the
parameter directions over the stack, the state form
of~\cite[Cond.~4]{gyorok2026obc}.
Gy\"or\"ok et al.\ note that an input designed with knowledge of the
structure of the omitted term can meet this condition~\cite[Remark~5]{gyorok2026obc}.
\todo{Missing: the instance. Whether the simulated truth meets $\Phi^\top\bm\delta=0$ belongs to Section~\ref{sec:setup}, which introduces it. Candidate for Section~\ref{sec:setup}: the Coulomb friction of the truth violates the condition, because its force has the sign of the velocity that multiplies the damping columns of \eqref{eq:obc_gantry_sens}, so OBC gives a unique split, not unbiased combinations (docs/references.md \texttt{gyorok2026obc} row; Section~VII bullet). The candidate must be checked against the normalised transition sensitivity, not only the force.}

The guarantee is narrower than the trained model.
The orthogonality of \eqref{eq:obc_construction} holds for one step on the
stack, not along trajectories or at the output.
Proposition~1 concerns the one-step criterion on the tuples.
$\mathcal M_{\mathrm{OBC}}$ minimises the closed-loop multistep criterion $V$
instead, so the proposition states what the correction permits, not the bias
of the trained model.
The tuples lie on $\bar x=0$.
The part of $f_{\mathrm{aug}}$ that acts through nonzero added states, which
PS2 activates, is therefore unconstrained.
$g_{\mathrm{aug}}$ is not corrected, because the baseline has no added-state
rows to be orthogonal to.
Orthogonality is Euclidean in the normalised coordinates of
\eqref{eq:aug_norm}, so the normalisation defines it.
The consistency and zero-covariance results
of~\cite[Thms.~16, 18]{gyorok2026obc} assume an input-output model linear in
its parameters and do not transfer.
